%%%%%%%%%%%%%%%%%%%%%%%%%%%%%%%%%%%%%%%%%%%%%%%%%%%%%%%%%%%%
% 12.7 TYPE THEORY
% The Composition of Advanced Mathematics
%%%%%%%%%%%%%%%%%%%%%%%%%%%%%%%%%%%%%%%%%%%%%%%%%%%%%%%%%%%%

\section{Type Theory}
\label{sec:type-theory}

\prerequisite{
Formal logic, proof theory, computability theory, recursion theory,
functions, relations, sets, mathematical induction, and basic programming
language concepts.
}

\learningobjectives{
By the end of this section, the reader should be able to:
\begin{itemize}
    \item explain the purpose of type theory;
    \item distinguish terms from types;
    \item understand typing judgments;
    \item recognize typing contexts;
    \item understand formation, introduction, elimination, and computation
          rules;
    \item understand simple types;
    \item recognize function types;
    \item understand product types;
    \item understand sum types;
    \item recognize unit and empty types;
    \item understand the Curry--Howard correspondence;
    \item interpret propositions as types;
    \item interpret proofs as terms;
    \item understand implication as function type;
    \item understand conjunction as product type;
    \item understand disjunction as sum type;
    \item understand falsity as an empty type;
    \item recognize beta reduction;
    \item recognize eta principles conceptually;
    \item understand normalization;
    \item recognize substitution in typed calculi;
    \item understand variable binding;
    \item recognize alpha conversion;
    \item understand simply typed lambda calculus;
    \item construct typed lambda terms;
    \item derive typing judgments;
    \item understand type safety conceptually;
    \item distinguish static and dynamic notions conceptually;
    \item understand dependent types;
    \item recognize dependent function types;
    \item recognize dependent pair types;
    \item understand universal quantification as dependent function type;
    \item understand existential quantification as dependent pair type;
    \item recognize identity types;
    \item understand equality proofs conceptually;
    \item recognize inductive types;
    \item understand natural numbers as an inductive type;
    \item understand recursion and induction principles;
    \item recognize lists and trees as inductive types;
    \item understand structural recursion;
    \item recognize propositions-as-types in constructive logic;
    \item understand proof-relevant mathematics;
    \item distinguish intensional and extensional equality conceptually;
    \item recognize universes;
    \item understand universe hierarchies;
    \item recognize dependent type theory as a foundational system;
    \item understand constructive content of proofs;
    \item recognize normalization as computation;
    \item understand canonicity conceptually;
    \item recognize decidability of type checking in suitable systems;
    \item understand definitional equality;
    \item distinguish definitional and propositional equality;
    \item recognize type inference conceptually;
    \item understand polymorphism conceptually;
    \item recognize parametric polymorphism;
    \item understand System F conceptually;
    \item recognize higher-order type systems;
    \item understand propositions-as-types in proof assistants;
    \item recognize proof terms and trusted kernels;
    \item understand propositions versus data types;
    \item recognize quotient-type difficulties conceptually;
    \item understand subtypes and refinement types conceptually;
    \item recognize homotopy type theory conceptually;
    \item understand univalence conceptually;
    \item recognize higher identity structure conceptually;
    \item connect type theory with logic, proof theory, programming
          languages, formal verification, constructive mathematics, and
          mathematical foundations.
\end{itemize}
}

%%%%%%%%%%%%%%%%%%%%%%%%%%%%%%%%%%%%%%%%%%%%%%%%%%%%%%%%%%%%
\subsection{What Is Type Theory?}
\label{subsec:what-is-type-theory}
%%%%%%%%%%%%%%%%%%%%%%%%%%%%%%%%%%%%%%%%%%%%%%%%%%%%%%%%%%%%

Type theory is a foundational framework in which mathematical objects are
organized according to their types.

Instead of beginning with one universal notion of set membership, type
theory begins with judgments such as

\[
\boxed{
a:A,
}
\]

read:

\[
\boxed{
\text{``\(a\) is a term of type \(A\).''}
}
\]

Types control how expressions may be formed and combined.

%%%%%%%%%%%%%%%%%%%%%%%%%%%%%%%%%%%%%%%%%%%%%%%%%%%%%%%%%%%%
\subsection{Terms and Types}
\label{subsec:terms-and-types}
%%%%%%%%%%%%%%%%%%%%%%%%%%%%%%%%%%%%%%%%%%%%%%%%%%%%%%%%%%%%

A \textbf{term} is an expression representing an object.

A \textbf{type} classifies terms.

For example,

\[
\boxed{
3:\mathbb{N}
}
\]

states that \(3\) is a natural number.

Likewise,

\[
\boxed{
f:A\to B
}
\]

states that \(f\) is a function taking inputs of type \(A\) and producing
outputs of type \(B\).

%%%%%%%%%%%%%%%%%%%%%%%%%%%%%%%%%%%%%%%%%%%%%%%%%%%%%%%%%%%%
\subsection{Typing Judgments}
\label{subsec:typing-judgments}
%%%%%%%%%%%%%%%%%%%%%%%%%%%%%%%%%%%%%%%%%%%%%%%%%%%%%%%%%%%%

The basic judgment

\[
\boxed{
a:A
}
\]

is different from an ordinary proposition about set membership.

It asserts that \(a\) is well formed as an object of type \(A\).

More generally, under assumptions \(\Gamma\), one writes

\[
\boxed{
\Gamma\vdash a:A.
}
\]

This means that term \(a\) has type \(A\) in context \(\Gamma\).

%%%%%%%%%%%%%%%%%%%%%%%%%%%%%%%%%%%%%%%%%%%%%%%%%%%%%%%%%%%%
\subsection{Contexts}
\label{subsec:type-theoretic-contexts}
%%%%%%%%%%%%%%%%%%%%%%%%%%%%%%%%%%%%%%%%%%%%%%%%%%%%%%%%%%%%

A context is a list of typed variable declarations.

For example,

\[
\boxed{
\Gamma
=
x:A,
y:B,
f:A\to B.
}
\]

A typing judgment such as

\[
\Gamma\vdash f(x):B
\]

is interpreted relative to those assumptions.

Contexts play a role analogous to local assumptions in logic.

%%%%%%%%%%%%%%%%%%%%%%%%%%%%%%%%%%%%%%%%%%%%%%%%%%%%%%%%%%%%
\subsection{Judgment Forms}
\label{subsec:type-theory-judgment-forms}
%%%%%%%%%%%%%%%%%%%%%%%%%%%%%%%%%%%%%%%%%%%%%%%%%%%%%%%%%%%%

Typical judgments include:

\[
\boxed{
A\ \text{type},
}
\]

\[
\boxed{
a:A,
}
\]

\[
\boxed{
A=B\ \text{type},
}
\]

and

\[
\boxed{
a=b:A.
}
\]

Different type theories formalize these judgments in different ways.

%%%%%%%%%%%%%%%%%%%%%%%%%%%%%%%%%%%%%%%%%%%%%%%%%%%%%%%%%%%%
\subsection{Four Kinds of Rules}
\label{subsec:type-theory-four-rules}
%%%%%%%%%%%%%%%%%%%%%%%%%%%%%%%%%%%%%%%%%%%%%%%%%%%%%%%%%%%%

A type constructor is often described by four kinds of rules:

\[
\boxed{
\text{formation},
}
\]

\[
\boxed{
\text{introduction},
}
\]

\[
\boxed{
\text{elimination},
}
\]

and

\[
\boxed{
\text{computation}.
}
\]

Formation says when a type exists.

Introduction says how to construct its terms.

Elimination says how to use those terms.

Computation says how introduction and elimination interact.

%%%%%%%%%%%%%%%%%%%%%%%%%%%%%%%%%%%%%%%%%%%%%%%%%%%%%%%%%%%%
\subsection{Function Types}
\label{subsec:type-theory-function-types}
%%%%%%%%%%%%%%%%%%%%%%%%%%%%%%%%%%%%%%%%%%%%%%%%%%%%%%%%%%%%

Given types \(A\) and \(B\), the function type is

\[
\boxed{
A\to B.
}
\]

A term

\[
f:A\to B
\]

takes an input

\[
a:A
\]

and produces

\[
\boxed{
f(a):B.
}
\]

%%%%%%%%%%%%%%%%%%%%%%%%%%%%%%%%%%%%%%%%%%%%%%%%%%%%%%%%%%%%
\subsection{Function Introduction}
\label{subsec:function-type-introduction}
%%%%%%%%%%%%%%%%%%%%%%%%%%%%%%%%%%%%%%%%%%%%%%%%%%%%%%%%%%%%

Suppose under assumption

\[
x:A
\]

we construct a term

\[
b:B.
\]

Then we obtain the function

\[
\boxed{
\lambda x.\,b
:
A\to B.
}
\]

The notation

\[
\lambda x.\,b
\]

represents a function taking \(x\) to \(b\).

%%%%%%%%%%%%%%%%%%%%%%%%%%%%%%%%%%%%%%%%%%%%%%%%%%%%%%%%%%%%
\subsection{Function Elimination}
\label{subsec:function-type-elimination}
%%%%%%%%%%%%%%%%%%%%%%%%%%%%%%%%%%%%%%%%%%%%%%%%%%%%%%%%%%%%

If

\[
f:A\to B
\]

and

\[
a:A,
\]

then function application gives

\[
\boxed{
f(a):B.
}
\]

Thus the elimination rule for a function type is ordinary function
application.

%%%%%%%%%%%%%%%%%%%%%%%%%%%%%%%%%%%%%%%%%%%%%%%%%%%%%%%%%%%%
\subsection{Beta Reduction}
\label{subsec:beta-reduction}
%%%%%%%%%%%%%%%%%%%%%%%%%%%%%%%%%%%%%%%%%%%%%%%%%%%%%%%%%%%%

Applying a lambda abstraction to an argument performs substitution.

The basic computation rule is

\[
\boxed{
(\lambda x.\,b)(a)
\longrightarrow
b[a/x].
}
\]

This is called beta reduction.

It is one of the fundamental computation rules of lambda calculus and type
theory.

%%%%%%%%%%%%%%%%%%%%%%%%%%%%%%%%%%%%%%%%%%%%%%%%%%%%%%%%%%%%
\subsection{Worked Example: Beta Reduction}
\label{subsec:worked-beta-reduction}
%%%%%%%%%%%%%%%%%%%%%%%%%%%%%%%%%%%%%%%%%%%%%%%%%%%%%%%%%%%%

\begin{workedexamplebox}{Evaluating a Lambda Term}

Consider

\[
(\lambda x.\,x+1)(4).
\]

Substitute

\[
4
\]

for \(x\):

\[
x+1
\longmapsto
4+1.
\]

Therefore,

\begin{answerbox}
\[
\boxed{
(\lambda x.\,x+1)(4)
\longrightarrow
5.
}
\]
\end{answerbox}

\end{workedexamplebox}

%%%%%%%%%%%%%%%%%%%%%%%%%%%%%%%%%%%%%%%%%%%%%%%%%%%%%%%%%%%%
\subsection{Simply Typed Lambda Calculus}
\label{subsec:simply-typed-lambda-calculus}
%%%%%%%%%%%%%%%%%%%%%%%%%%%%%%%%%%%%%%%%%%%%%%%%%%%%%%%%%%%%

The simply typed lambda calculus begins with basic types and constructs new
types using function types.

Typical terms include:

\[
\boxed{
x,
}
\]

\[
\boxed{
\lambda x.\,t,
}
\]

and

\[
\boxed{
t\,u.
}
\]

Typing rules prevent meaningless applications.

%%%%%%%%%%%%%%%%%%%%%%%%%%%%%%%%%%%%%%%%%%%%%%%%%%%%%%%%%%%%
\subsection{Function Typing Rule}
\label{subsec:function-typing-rule}
%%%%%%%%%%%%%%%%%%%%%%%%%%%%%%%%%%%%%%%%%%%%%%%%%%%%%%%%%%%%

Lambda abstraction has typing rule

\[
\boxed{
\frac{
\Gamma,x:A
\vdash
b:B
}{
\Gamma
\vdash
\lambda x.\,b
:
A\to B
}.
}
\]

Application has rule

\[
\boxed{
\frac{
\Gamma\vdash f:A\to B
\qquad
\Gamma\vdash a:A
}{
\Gamma\vdash f(a):B
}.
}
\]

%%%%%%%%%%%%%%%%%%%%%%%%%%%%%%%%%%%%%%%%%%%%%%%%%%%%%%%%%%%%
\subsection{Worked Example: Typed Identity Function}
\label{subsec:worked-typed-identity}
%%%%%%%%%%%%%%%%%%%%%%%%%%%%%%%%%%%%%%%%%%%%%%%%%%%%%%%%%%%%

\begin{workedexamplebox}{Identity Function}

Assume

\[
x:A.
\]

Then trivially,

\[
x:A.
\]

By function introduction,

\[
\lambda x.\,x
:
A\to A.
\]

Thus

\begin{answerbox}
\[
\boxed{
\lambda x.\,x
:
A\to A.
}
\]
\end{answerbox}

This term represents the identity function on type \(A\).

\end{workedexamplebox}

%%%%%%%%%%%%%%%%%%%%%%%%%%%%%%%%%%%%%%%%%%%%%%%%%%%%%%%%%%%%
\subsection{Product Types}
\label{subsec:product-types}
%%%%%%%%%%%%%%%%%%%%%%%%%%%%%%%%%%%%%%%%%%%%%%%%%%%%%%%%%%%%

Given types \(A\) and \(B\), their product type is written

\[
\boxed{
A\times B.
}
\]

A term of this type is a pair

\[
\boxed{
(a,b)
}
\]

with

\[
a:A
\]

and

\[
b:B.
\]

%%%%%%%%%%%%%%%%%%%%%%%%%%%%%%%%%%%%%%%%%%%%%%%%%%%%%%%%%%%%
\subsection{Product Introduction}
\label{subsec:product-introduction}
%%%%%%%%%%%%%%%%%%%%%%%%%%%%%%%%%%%%%%%%%%%%%%%%%%%%%%%%%%%%

The introduction rule is

\[
\boxed{
\frac{
a:A
\qquad
b:B
}{
(a,b):A\times B
}.
}
\]

Thus a term of a product type contains both pieces of data.

%%%%%%%%%%%%%%%%%%%%%%%%%%%%%%%%%%%%%%%%%%%%%%%%%%%%%%%%%%%%
\subsection{Product Elimination}
\label{subsec:product-elimination}
%%%%%%%%%%%%%%%%%%%%%%%%%%%%%%%%%%%%%%%%%%%%%%%%%%%%%%%%%%%%

Given

\[
p:A\times B,
\]

one may extract the two components:

\[
\boxed{
\pi_1(p):A,
}
\]

and

\[
\boxed{
\pi_2(p):B.
}
\]

For a pair,

\[
\boxed{
\pi_1(a,b)=a,
}
\]

\[
\boxed{
\pi_2(a,b)=b.
}
\]

%%%%%%%%%%%%%%%%%%%%%%%%%%%%%%%%%%%%%%%%%%%%%%%%%%%%%%%%%%%%
\subsection{Sum Types}
\label{subsec:sum-types}
%%%%%%%%%%%%%%%%%%%%%%%%%%%%%%%%%%%%%%%%%%%%%%%%%%%%%%%%%%%%

Given types \(A\) and \(B\), their sum type is

\[
\boxed{
A+B.
}
\]

A term of \(A+B\) is either:

\[
\boxed{
\operatorname{inl}(a)
}
\]

for

\[
a:A,
\]

or

\[
\boxed{
\operatorname{inr}(b)
}
\]

for

\[
b:B.
\]

The tags identify which side was used.

%%%%%%%%%%%%%%%%%%%%%%%%%%%%%%%%%%%%%%%%%%%%%%%%%%%%%%%%%%%%
\subsection{Sum Elimination}
\label{subsec:sum-elimination}
%%%%%%%%%%%%%%%%%%%%%%%%%%%%%%%%%%%%%%%%%%%%%%%%%%%%%%%%%%%%

To define a result from

\[
z:A+B,
\]

one provides:

\[
\boxed{
\text{a rule for the \(A\)-case},
}
\]

and

\[
\boxed{
\text{a rule for the \(B\)-case}.
}
\]

This is typed case analysis.

%%%%%%%%%%%%%%%%%%%%%%%%%%%%%%%%%%%%%%%%%%%%%%%%%%%%%%%%%%%%
\subsection{Unit Type}
\label{subsec:unit-type}
%%%%%%%%%%%%%%%%%%%%%%%%%%%%%%%%%%%%%%%%%%%%%%%%%%%%%%%%%%%%

The unit type, often written

\[
\boxed{
\mathbf{1},
}
\]

has exactly one canonical term, commonly written

\[
\boxed{
\star:\mathbf{1}.
}
\]

It represents a type containing no significant information beyond
inhabitation.

%%%%%%%%%%%%%%%%%%%%%%%%%%%%%%%%%%%%%%%%%%%%%%%%%%%%%%%%%%%%
\subsection{Empty Type}
\label{subsec:empty-type}
%%%%%%%%%%%%%%%%%%%%%%%%%%%%%%%%%%%%%%%%%%%%%%%%%%%%%%%%%%%%

The empty type, often written

\[
\boxed{
\mathbf{0},
}
\]

has no terms.

Thus there is no term

\[
\boxed{
a:\mathbf{0}.
}
\]

From an impossible term of \(\mathbf{0}\), one can derive a term of any
type.

This corresponds logically to explosion from contradiction.

%%%%%%%%%%%%%%%%%%%%%%%%%%%%%%%%%%%%%%%%%%%%%%%%%%%%%%%%%%%%
\subsection{Curry--Howard Correspondence}
\label{subsec:type-theory-curry-howard}
%%%%%%%%%%%%%%%%%%%%%%%%%%%%%%%%%%%%%%%%%%%%%%%%%%%%%%%%%%%%

The Curry--Howard correspondence connects logic and type theory:

\[
\boxed{
\text{propositions}
\leftrightarrow
\text{types},
}
\]

\[
\boxed{
\text{proofs}
\leftrightarrow
\text{terms}.
}
\]

Under this interpretation,

\[
\boxed{
\text{proving proposition }P
}
\]

means

\[
\boxed{
\text{constructing a term }p:P.
}
\]

%%%%%%%%%%%%%%%%%%%%%%%%%%%%%%%%%%%%%%%%%%%%%%%%%%%%%%%%%%%%
\subsection{Implication as Function Type}
\label{subsec:type-theory-implication}
%%%%%%%%%%%%%%%%%%%%%%%%%%%%%%%%%%%%%%%%%%%%%%%%%%%%%%%%%%%%

A proof of

\[
P\rightarrow Q
\]

takes a proof of \(P\) and transforms it into a proof of \(Q\).

Thus

\[
\boxed{
P\rightarrow Q
}
\]

corresponds directly to the function type

\[
\boxed{
P\to Q.
}
\]

A proof is a function from evidence of \(P\) to evidence of \(Q\).

%%%%%%%%%%%%%%%%%%%%%%%%%%%%%%%%%%%%%%%%%%%%%%%%%%%%%%%%%%%%
\subsection{Conjunction as Product Type}
\label{subsec:type-theory-conjunction}
%%%%%%%%%%%%%%%%%%%%%%%%%%%%%%%%%%%%%%%%%%%%%%%%%%%%%%%%%%%%

A proof of

\[
P\land Q
\]

contains:

\[
\boxed{
p:P
}
\]

and

\[
\boxed{
q:Q.
}
\]

Thus conjunction corresponds to

\[
\boxed{
P\times Q.
}
\]

The proof term is

\[
\boxed{
(p,q):P\times Q.
}
\]

%%%%%%%%%%%%%%%%%%%%%%%%%%%%%%%%%%%%%%%%%%%%%%%%%%%%%%%%%%%%
\subsection{Disjunction as Sum Type}
\label{subsec:type-theory-disjunction}
%%%%%%%%%%%%%%%%%%%%%%%%%%%%%%%%%%%%%%%%%%%%%%%%%%%%%%%%%%%%

A constructive proof of

\[
P\lor Q
\]

must identify which alternative holds.

Thus disjunction corresponds to

\[
\boxed{
P+Q.
}
\]

A proof is either

\[
\boxed{
\operatorname{inl}(p)
}
\]

or

\[
\boxed{
\operatorname{inr}(q).
}
\]

%%%%%%%%%%%%%%%%%%%%%%%%%%%%%%%%%%%%%%%%%%%%%%%%%%%%%%%%%%%%
\subsection{Falsity as Empty Type}
\label{subsec:falsity-empty-type}
%%%%%%%%%%%%%%%%%%%%%%%%%%%%%%%%%%%%%%%%%%%%%%%%%%%%%%%%%%%%

Falsity corresponds to the empty type:

\[
\boxed{
\bot
\leftrightarrow
\mathbf{0}.
}
\]

A proof of falsity would be a term

\[
\boxed{
p:\mathbf{0}.
}
\]

Because no such canonical term exists, the type represents impossibility.

%%%%%%%%%%%%%%%%%%%%%%%%%%%%%%%%%%%%%%%%%%%%%%%%%%%%%%%%%%%%
\subsection{Negation as Function Type}
\label{subsec:type-theory-negation}
%%%%%%%%%%%%%%%%%%%%%%%%%%%%%%%%%%%%%%%%%%%%%%%%%%%%%%%%%%%%

Negation is interpreted as

\[
\boxed{
\neg P
=
P\to\mathbf{0}.
}
\]

Thus a proof of \(\neg P\) is a function that turns any proof of \(P\) into
a contradiction.

%%%%%%%%%%%%%%%%%%%%%%%%%%%%%%%%%%%%%%%%%%%%%%%%%%%%%%%%%%%%
\subsection{Worked Example: Proof as Program}
\label{subsec:worked-proof-as-program}
%%%%%%%%%%%%%%%%%%%%%%%%%%%%%%%%%%%%%%%%%%%%%%%%%%%%%%%%%%%%

\begin{workedexamplebox}{Proof of \(P\land Q\rightarrow P\)}

Suppose

\[
z:P\times Q.
\]

The first projection gives

\[
\pi_1(z):P.
\]

Therefore the function

\[
\lambda z.\,\pi_1(z)
\]

has type

\[
P\times Q\to P.
\]

\begin{answerbox}
\[
\boxed{
\lambda z.\,\pi_1(z)
:
P\times Q\to P.
}
\]
\end{answerbox}

Under Curry--Howard, this term is a proof of

\[
P\land Q\rightarrow P.
\]

\end{workedexamplebox}

%%%%%%%%%%%%%%%%%%%%%%%%%%%%%%%%%%%%%%%%%%%%%%%%%%%%%%%%%%%%
\subsection{Alpha Conversion}
\label{subsec:alpha-conversion-type-theory}
%%%%%%%%%%%%%%%%%%%%%%%%%%%%%%%%%%%%%%%%%%%%%%%%%%%%%%%%%%%%

Bound variable names do not affect the essential meaning of a lambda term.

Thus

\[
\boxed{
\lambda x.\,x
}
\]

and

\[
\boxed{
\lambda y.\,y
}
\]

represent the same binding structure.

Renaming bound variables is called alpha conversion.

%%%%%%%%%%%%%%%%%%%%%%%%%%%%%%%%%%%%%%%%%%%%%%%%%%%%%%%%%%%%
\subsection{Capture-Avoiding Substitution}
\label{subsec:type-theory-substitution}
%%%%%%%%%%%%%%%%%%%%%%%%%%%%%%%%%%%%%%%%%%%%%%%%%%%%%%%%%%%%

Substitution must avoid accidentally binding free variables.

For example, substituting \(y\) into

\[
\lambda y.\,x
\]

for \(x\) requires renaming the bound variable first.

One may rewrite

\[
\lambda y.\,x
\]

as

\[
\lambda z.\,x
\]

and then substitute:

\[
\boxed{
(\lambda z.\,x)[y/x]
=
\lambda z.\,y.
}
\]

%%%%%%%%%%%%%%%%%%%%%%%%%%%%%%%%%%%%%%%%%%%%%%%%%%%%%%%%%%%%
\subsection{Eta Principle}
\label{subsec:eta-principle}
%%%%%%%%%%%%%%%%%%%%%%%%%%%%%%%%%%%%%%%%%%%%%%%%%%%%%%%%%%%%

A function \(f:A\to B\) is often considered extensionally equivalent to

\[
\boxed{
\lambda x.\,f(x).
}
\]

This is associated with the eta principle.

Conceptually,

\[
\boxed{
\lambda x.\,f(x)
\approx
f
}
\]

when \(x\) does not occur freely in \(f\).

Exact equality conventions depend on the type theory.

%%%%%%%%%%%%%%%%%%%%%%%%%%%%%%%%%%%%%%%%%%%%%%%%%%%%%%%%%%%%
\subsection{Normalization}
\label{subsec:type-theory-normalization}
%%%%%%%%%%%%%%%%%%%%%%%%%%%%%%%%%%%%%%%%%%%%%%%%%%%%%%%%%%%%

Normalization repeatedly performs computation rules such as beta
reduction.

A term is in normal form when no permitted reduction step remains.

For example,

\[
(\lambda x.\,x)(a)
\]

reduces to

\[
\boxed{
a.
}
\]

Normalization connects proof simplification with program evaluation.

%%%%%%%%%%%%%%%%%%%%%%%%%%%%%%%%%%%%%%%%%%%%%%%%%%%%%%%%%%%%
\subsection{Strong Normalization}
\label{subsec:type-theory-strong-normalization}
%%%%%%%%%%%%%%%%%%%%%%%%%%%%%%%%%%%%%%%%%%%%%%%%%%%%%%%%%%%%

A calculus is strongly normalizing if every valid reduction sequence
eventually terminates.

The simply typed lambda calculus is strongly normalizing.

Thus well-typed simply typed lambda terms cannot reduce forever.

This is a major structural property.

%%%%%%%%%%%%%%%%%%%%%%%%%%%%%%%%%%%%%%%%%%%%%%%%%%%%%%%%%%%%
\subsection{Type Safety}
\label{subsec:type-safety}
%%%%%%%%%%%%%%%%%%%%%%%%%%%%%%%%%%%%%%%%%%%%%%%%%%%%%%%%%%%%

Type safety means, informally, that well-typed terms do not encounter
certain kinds of meaningless computation.

Two important metatheoretic ideas are:

\[
\boxed{
\text{preservation},
}
\]

and

\[
\boxed{
\text{progress}.
}
\]

%%%%%%%%%%%%%%%%%%%%%%%%%%%%%%%%%%%%%%%%%%%%%%%%%%%%%%%%%%%%
\subsection{Preservation}
\label{subsec:type-preservation}
%%%%%%%%%%%%%%%%%%%%%%%%%%%%%%%%%%%%%%%%%%%%%%%%%%%%%%%%%%%%

Preservation states that reduction does not change a term's type.

If

\[
\boxed{
t:A
}
\]

and

\[
\boxed{
t\longrightarrow t',
}
\]

then one expects

\[
\boxed{
t':A.
}
\]

This is also called subject reduction.

%%%%%%%%%%%%%%%%%%%%%%%%%%%%%%%%%%%%%%%%%%%%%%%%%%%%%%%%%%%%
\subsection{Progress}
\label{subsec:type-progress}
%%%%%%%%%%%%%%%%%%%%%%%%%%%%%%%%%%%%%%%%%%%%%%%%%%%%%%%%%%%%

Progress states, in a suitable closed typed language, that a well-typed term
is either:

\[
\boxed{
\text{already a value},
}
\]

or

\[
\boxed{
\text{can take another computation step}.
}
\]

It therefore does not become stuck because of a type mismatch.

%%%%%%%%%%%%%%%%%%%%%%%%%%%%%%%%%%%%%%%%%%%%%%%%%%%%%%%%%%%%
\subsection{Dependent Types}
\label{subsec:dependent-types}
%%%%%%%%%%%%%%%%%%%%%%%%%%%%%%%%%%%%%%%%%%%%%%%%%%%%%%%%%%%%

In simple type theory, the output type of a function does not depend on the
specific input value.

Dependent type theory allows types to depend on terms.

Thus one may have a family of types

\[
\boxed{
B(x)
}
\]

indexed by

\[
x:A.
\]

This greatly increases expressive power.

%%%%%%%%%%%%%%%%%%%%%%%%%%%%%%%%%%%%%%%%%%%%%%%%%%%%%%%%%%%%
\subsection{Dependent Function Types}
\label{subsec:dependent-function-types}
%%%%%%%%%%%%%%%%%%%%%%%%%%%%%%%%%%%%%%%%%%%%%%%%%%%%%%%%%%%%

The dependent function type is written

\[
\boxed{
\prod_{x:A}B(x).
}
\]

It is also called a \(\Pi\)-type.

A term

\[
f:
\prod_{x:A}B(x)
\]

assigns to each

\[
x:A
\]

a term

\[
\boxed{
f(x):B(x).
}
\]

The codomain type may vary with \(x\).

%%%%%%%%%%%%%%%%%%%%%%%%%%%%%%%%%%%%%%%%%%%%%%%%%%%%%%%%%%%%
\subsection{Ordinary Functions as Special Pi Types}
\label{subsec:ordinary-function-pi-type}
%%%%%%%%%%%%%%%%%%%%%%%%%%%%%%%%%%%%%%%%%%%%%%%%%%%%%%%%%%%%

If

\[
B(x)=B
\]

does not actually depend on \(x\), then

\[
\boxed{
\prod_{x:A}B
}
\]

is simply the ordinary function type

\[
\boxed{
A\to B.
}
\]

Thus ordinary function types are special cases of dependent function types.

%%%%%%%%%%%%%%%%%%%%%%%%%%%%%%%%%%%%%%%%%%%%%%%%%%%%%%%%%%%%
\subsection{Dependent Pair Types}
\label{subsec:dependent-pair-types}
%%%%%%%%%%%%%%%%%%%%%%%%%%%%%%%%%%%%%%%%%%%%%%%%%%%%%%%%%%%%

The dependent pair type is written

\[
\boxed{
\sum_{x:A}B(x).
}
\]

It is also called a \(\Sigma\)-type.

A term has form

\[
\boxed{
(a,b)
}
\]

where

\[
a:A
\]

and

\[
b:B(a).
\]

The type of the second component depends on the first.

%%%%%%%%%%%%%%%%%%%%%%%%%%%%%%%%%%%%%%%%%%%%%%%%%%%%%%%%%%%%
\subsection{Ordinary Products as Special Sigma Types}
\label{subsec:ordinary-product-sigma-type}
%%%%%%%%%%%%%%%%%%%%%%%%%%%%%%%%%%%%%%%%%%%%%%%%%%%%%%%%%%%%

If \(B(x)=B\) is constant, then

\[
\boxed{
\sum_{x:A}B
}
\]

behaves like

\[
\boxed{
A\times B.
}
\]

Thus ordinary products are special cases of dependent pair types.

%%%%%%%%%%%%%%%%%%%%%%%%%%%%%%%%%%%%%%%%%%%%%%%%%%%%%%%%%%%%
\subsection{Universal Quantification as Pi Type}
\label{subsec:forall-pi-type}
%%%%%%%%%%%%%%%%%%%%%%%%%%%%%%%%%%%%%%%%%%%%%%%%%%%%%%%%%%%%

Under Curry--Howard, the proposition

\[
\forall x:A\,P(x)
\]

corresponds to

\[
\boxed{
\prod_{x:A}P(x).
}
\]

A proof must assign to every

\[
a:A
\]

a proof

\[
\boxed{
p(a):P(a).
}
\]

Thus universal proof is a dependent function.

%%%%%%%%%%%%%%%%%%%%%%%%%%%%%%%%%%%%%%%%%%%%%%%%%%%%%%%%%%%%
\subsection{Existential Quantification as Sigma Type}
\label{subsec:exists-sigma-type}
%%%%%%%%%%%%%%%%%%%%%%%%%%%%%%%%%%%%%%%%%%%%%%%%%%%%%%%%%%%%

The proposition

\[
\exists x:A\,P(x)
\]

corresponds constructively to

\[
\boxed{
\sum_{x:A}P(x).
}
\]

A proof contains:

\[
\boxed{
\text{a witness }a:A
}
\]

and

\[
\boxed{
\text{a proof }p:P(a).
}
\]

Thus existential proof contains computationally meaningful data.

%%%%%%%%%%%%%%%%%%%%%%%%%%%%%%%%%%%%%%%%%%%%%%%%%%%%%%%%%%%%
\subsection{Worked Example: Existential Proof}
\label{subsec:worked-existential-sigma}
%%%%%%%%%%%%%%%%%%%%%%%%%%%%%%%%%%%%%%%%%%%%%%%%%%%%%%%%%%%%

\begin{workedexamplebox}{Constructing a Witness}

Consider the proposition

\[
\exists n:\mathbb{N}\,
(n=4).
\]

A constructive proof supplies the witness

\[
4:\mathbb{N}
\]

together with evidence

\[
p:4=4.
\]

Thus the proof is represented by the dependent pair

\[
\begin{answerbox}
\[
\boxed{
(4,p)
:
\sum_{n:\mathbb{N}}
(n=4).
}
\]
\end{answerbox}

\end{workedexamplebox}

%%%%%%%%%%%%%%%%%%%%%%%%%%%%%%%%%%%%%%%%%%%%%%%%%%%%%%%%%%%%
\subsection{Identity Types}
\label{subsec:identity-types}
%%%%%%%%%%%%%%%%%%%%%%%%%%%%%%%%%%%%%%%%%%%%%%%%%%%%%%%%%%%%

For terms

\[
a,b:A,
\]

dependent type theory may contain an identity type

\[
\boxed{
a=_A b.
}
\]

A term

\[
p:a=_A b
\]

is evidence that \(a\) and \(b\) are equal.

Thus equality itself becomes a type with possible proof terms.

%%%%%%%%%%%%%%%%%%%%%%%%%%%%%%%%%%%%%%%%%%%%%%%%%%%%%%%%%%%%
\subsection{Reflexivity}
\label{subsec:identity-reflexivity}
%%%%%%%%%%%%%%%%%%%%%%%%%%%%%%%%%%%%%%%%%%%%%%%%%%%%%%%%%%%%

Every term is equal to itself.

Identity introduction therefore provides

\[
\boxed{
\operatorname{refl}_a
:
a=_A a.
}
\]

This is the canonical constructor for identity types.

%%%%%%%%%%%%%%%%%%%%%%%%%%%%%%%%%%%%%%%%%%%%%%%%%%%%%%%%%%%%
\subsection{Equality Elimination}
\label{subsec:identity-elimination}
%%%%%%%%%%%%%%%%%%%%%%%%%%%%%%%%%%%%%%%%%%%%%%%%%%%%%%%%%%%%

The elimination principle for identity types says that proofs involving
arbitrary equality evidence can be reduced to reasoning about reflexivity.

This principle is often called

\[
\boxed{
J
}
\]

or path induction in homotopy-inspired terminology.

It allows properties to be transported along equalities.

%%%%%%%%%%%%%%%%%%%%%%%%%%%%%%%%%%%%%%%%%%%%%%%%%%%%%%%%%%%%
\subsection{Transport}
\label{subsec:type-theory-transport}
%%%%%%%%%%%%%%%%%%%%%%%%%%%%%%%%%%%%%%%%%%%%%%%%%%%%%%%%%%%%

Suppose

\[
P:A\to\mathcal{U}
\]

is a family of types and

\[
p:a=_A b.
\]

Then a term

\[
u:P(a)
\]

can be transported along \(p\) to obtain a term of

\[
\boxed{
P(b).
}
\]

Equality therefore supports substitution inside dependent types.

%%%%%%%%%%%%%%%%%%%%%%%%%%%%%%%%%%%%%%%%%%%%%%%%%%%%%%%%%%%%
\subsection{Definitional Equality}
\label{subsec:definitional-equality}
%%%%%%%%%%%%%%%%%%%%%%%%%%%%%%%%%%%%%%%%%%%%%%%%%%%%%%%%%%%%

Definitional equality means two expressions compute or simplify to the same
term according to the formal reduction rules.

For example,

\[
\boxed{
(\lambda x.\,x)(a)
\equiv
a
}
\]

by beta computation.

This equality is usually checked by computation rather than by supplying a
proof term.

%%%%%%%%%%%%%%%%%%%%%%%%%%%%%%%%%%%%%%%%%%%%%%%%%%%%%%%%%%%%
\subsection{Propositional Equality}
\label{subsec:propositional-equality}
%%%%%%%%%%%%%%%%%%%%%%%%%%%%%%%%%%%%%%%%%%%%%%%%%%%%%%%%%%%%

Propositional equality is represented by an identity type.

To establish

\[
a=_A b,
\]

one provides a term

\[
\boxed{
p:a=_A b.
}
\]

Thus definitional equality and propositional equality play different roles.

%%%%%%%%%%%%%%%%%%%%%%%%%%%%%%%%%%%%%%%%%%%%%%%%%%%%%%%%%%%%
\subsection{Worked Example: Two Equalities}
\label{subsec:worked-two-equalities}
%%%%%%%%%%%%%%%%%%%%%%%%%%%%%%%%%%%%%%%%%%%%%%%%%%%%%%%%%%%%

\begin{workedexamplebox}{Computation Versus Proof}

The expression

\[
(\lambda x.\,x)(3)
\]

reduces directly to

\[
3.
\]

This is definitional equality.

By contrast, a theorem asserting equality between two mathematically
derived expressions may require an explicit proof term.

\begin{answerbox}
\[
\boxed{
\text{definitional equality}
=
\text{computation},
}
\]

\[
\boxed{
\text{propositional equality}
=
\text{inhabited identity type}.
}
\]
\end{answerbox}

\end{workedexamplebox}

%%%%%%%%%%%%%%%%%%%%%%%%%%%%%%%%%%%%%%%%%%%%%%%%%%%%%%%%%%%%
\subsection{Inductive Types}
\label{subsec:inductive-types}
%%%%%%%%%%%%%%%%%%%%%%%%%%%%%%%%%%%%%%%%%%%%%%%%%%%%%%%%%%%%

An inductive type is specified by constructors that generate its canonical
terms.

The natural numbers provide the basic example.

One declares constructors

\[
\boxed{
0:\mathbb{N}
}
\]

and

\[
\boxed{
S:\mathbb{N}\to\mathbb{N}.
}
\]

Every canonical natural number is generated from these constructors.

%%%%%%%%%%%%%%%%%%%%%%%%%%%%%%%%%%%%%%%%%%%%%%%%%%%%%%%%%%%%
\subsection{Natural Number Recursion}
\label{subsec:natural-number-recursion-type-theory}
%%%%%%%%%%%%%%%%%%%%%%%%%%%%%%%%%%%%%%%%%%%%%%%%%%%%%%%%%%%%

To define

\[
f:\mathbb{N}\to C,
\]

one provides:

\[
\boxed{
c_0:C,
}
\]

and

\[
\boxed{
s:C\to C
}
\]

in the simplest constant-family recursion setting.

Then:

\[
\boxed{
f(0)=c_0,
}
\]

\[
\boxed{
f(Sn)=s(f(n)).
}
\]

More general recursors may depend explicitly on \(n\).

%%%%%%%%%%%%%%%%%%%%%%%%%%%%%%%%%%%%%%%%%%%%%%%%%%%%%%%%%%%%
\subsection{Natural Number Induction}
\label{subsec:natural-number-induction-type-theory}
%%%%%%%%%%%%%%%%%%%%%%%%%%%%%%%%%%%%%%%%%%%%%%%%%%%%%%%%%%%%

Let

\[
P:\mathbb{N}\to\mathcal{U}
\]

be a type family.

To construct

\[
\boxed{
\prod_{n:\mathbb{N}}P(n),
}
\]

it is enough to construct:

\[
\boxed{
p_0:P(0),
}
\]

and

\[
\boxed{
s:
\prod_{n:\mathbb{N}}
P(n)\to P(Sn).
}
\]

This is the dependent induction principle for natural numbers.

%%%%%%%%%%%%%%%%%%%%%%%%%%%%%%%%%%%%%%%%%%%%%%%%%%%%%%%%%%%%
\subsection{Recursion Versus Induction}
\label{subsec:type-theory-recursion-vs-induction}
%%%%%%%%%%%%%%%%%%%%%%%%%%%%%%%%%%%%%%%%%%%%%%%%%%%%%%%%%%%%

Recursion defines data or functions.

Induction proves dependent properties.

Recursion has a fixed output type \(C\).

Induction allows the target type

\[
\boxed{
P(n)
}
\]

to vary with the natural number \(n\).

Thus induction is a dependent generalization of recursion.

%%%%%%%%%%%%%%%%%%%%%%%%%%%%%%%%%%%%%%%%%%%%%%%%%%%%%%%%%%%%
\subsection{Lists}
\label{subsec:type-theory-lists}
%%%%%%%%%%%%%%%%%%%%%%%%%%%%%%%%%%%%%%%%%%%%%%%%%%%%%%%%%%%%

For a type \(A\), the list type

\[
\boxed{
\operatorname{List}(A)
}
\]

may be generated by constructors

\[
\boxed{
\operatorname{nil}
:
\operatorname{List}(A),
}
\]

and

\[
\boxed{
\operatorname{cons}
:
A
\to
\operatorname{List}(A)
\to
\operatorname{List}(A).
}
\]

These constructors determine structural recursion on lists.

%%%%%%%%%%%%%%%%%%%%%%%%%%%%%%%%%%%%%%%%%%%%%%%%%%%%%%%%%%%%
\subsection{Structural Recursion}
\label{subsec:structural-recursion-type-theory}
%%%%%%%%%%%%%%%%%%%%%%%%%%%%%%%%%%%%%%%%%%%%%%%%%%%%%%%%%%%%

A function on an inductively generated object can often be defined by
specifying its behavior on each constructor.

For lists:

\[
f(\operatorname{nil})
=
c,
\]

and

\[
f(\operatorname{cons}(a,\ell))
=
F(a,\ell,f(\ell)).
\]

Because recursive calls are made on structurally smaller data, termination
is built into the definition.

%%%%%%%%%%%%%%%%%%%%%%%%%%%%%%%%%%%%%%%%%%%%%%%%%%%%%%%%%%%%
\subsection{Binary Trees}
\label{subsec:type-theory-binary-trees}
%%%%%%%%%%%%%%%%%%%%%%%%%%%%%%%%%%%%%%%%%%%%%%%%%%%%%%%%%%%%

A binary tree type may be generated by constructors such as

\[
\boxed{
\operatorname{leaf}:A\to\operatorname{Tree}(A),
}
\]

and

\[
\boxed{
\operatorname{node}:
\operatorname{Tree}(A)
\to
\operatorname{Tree}(A)
\to
\operatorname{Tree}(A).
}
\]

Functions over trees are defined recursively according to this shape.

%%%%%%%%%%%%%%%%%%%%%%%%%%%%%%%%%%%%%%%%%%%%%%%%%%%%%%%%%%%%
\subsection{Inductive Definitions and Termination}
\label{subsec:inductive-definition-termination}
%%%%%%%%%%%%%%%%%%%%%%%%%%%%%%%%%%%%%%%%%%%%%%%%%%%%%%%%%%%%

Structural recursion on strictly smaller constructor components provides an
intrinsic termination argument.

Thus the type system may reject recursive definitions that cannot be shown
to decrease structurally.

This is important in systems where every accepted function must be total.

%%%%%%%%%%%%%%%%%%%%%%%%%%%%%%%%%%%%%%%%%%%%%%%%%%%%%%%%%%%%
\subsection{Canonicity}
\label{subsec:type-theory-canonicity}
%%%%%%%%%%%%%%%%%%%%%%%%%%%%%%%%%%%%%%%%%%%%%%%%%%%%%%%%%%%%

Canonicity states, roughly, that every closed term of a basic data type
computes to a canonical constructor form.

For natural numbers, one expects every closed term

\[
n:\mathbb{N}
\]

to reduce to

\[
\boxed{
0,
\quad
S0,
\quad
SS0,
\quad
\ldots.
}
\]

This connects formal proof with effective computation.

%%%%%%%%%%%%%%%%%%%%%%%%%%%%%%%%%%%%%%%%%%%%%%%%%%%%%%%%%%%%
\subsection{Normalization and Consistency}
\label{subsec:type-theory-normalization-consistency}
%%%%%%%%%%%%%%%%%%%%%%%%%%%%%%%%%%%%%%%%%%%%%%%%%%%%%%%%%%%%

Suppose every closed well-typed term normalizes.

If there were a closed term

\[
p:\mathbf{0},
\]

it would normalize to a canonical term of the empty type.

But the empty type has no constructors.

Therefore, normalization and canonicity can support consistency proofs for
appropriate type theories.

%%%%%%%%%%%%%%%%%%%%%%%%%%%%%%%%%%%%%%%%%%%%%%%%%%%%%%%%%%%%
\subsection{Universes}
\label{subsec:type-theory-universes}
%%%%%%%%%%%%%%%%%%%%%%%%%%%%%%%%%%%%%%%%%%%%%%%%%%%%%%%%%%%%

A universe is a type whose terms represent types.

One may write

\[
\boxed{
A:\mathcal{U}.
}
\]

This means \(A\) is a small type belonging to universe \(\mathcal{U}\).

Universes permit type theory to reason internally about collections of
types.

%%%%%%%%%%%%%%%%%%%%%%%%%%%%%%%%%%%%%%%%%%%%%%%%%%%%%%%%%%%%
\subsection{Universe Hierarchy}
\label{subsec:universe-hierarchy}
%%%%%%%%%%%%%%%%%%%%%%%%%%%%%%%%%%%%%%%%%%%%%%%%%%%%%%%%%%%%

Naively assuming

\[
\mathcal{U}:\mathcal{U}
\]

can lead to paradoxes in sufficiently expressive systems.

A common solution uses a hierarchy

\[
\boxed{
\mathcal{U}_0,
\mathcal{U}_1,
\mathcal{U}_2,
\ldots
}
\]

with

\[
\boxed{
\mathcal{U}_0:\mathcal{U}_1,
\qquad
\mathcal{U}_1:\mathcal{U}_2,
\quad\ldots.
}
\]

This resembles size hierarchies in set theory.

%%%%%%%%%%%%%%%%%%%%%%%%%%%%%%%%%%%%%%%%%%%%%%%%%%%%%%%%%%%%
\subsection{Type-Theoretic Paradoxes}
\label{subsec:type-theoretic-paradoxes}
%%%%%%%%%%%%%%%%%%%%%%%%%%%%%%%%%%%%%%%%%%%%%%%%%%%%%%%%%%%%

Just as unrestricted set formation produces Russell's paradox, unrestricted
self-referential type formation can create contradictions.

Modern type theories carefully restrict universe formation and dependency
to maintain consistency.

Thus foundational design requires balancing:

\[
\boxed{
\text{expressive power}
}
\]

and

\[
\boxed{
\text{logical safety}.
}
\]

%%%%%%%%%%%%%%%%%%%%%%%%%%%%%%%%%%%%%%%%%%%%%%%%%%%%%%%%%%%%
\subsection{Intensional Type Theory}
\label{subsec:intensional-type-theory}
%%%%%%%%%%%%%%%%%%%%%%%%%%%%%%%%%%%%%%%%%%%%%%%%%%%%%%%%%%%%

In intensional type theory, equality proofs carry structure and
propositional equality is not automatically identified with definitional
equality.

Thus two terms may be propositionally equal without reducing to the same
syntactic normal form.

This distinction supports computational interpretation.

%%%%%%%%%%%%%%%%%%%%%%%%%%%%%%%%%%%%%%%%%%%%%%%%%%%%%%%%%%%%
\subsection{Extensional Type Theory}
\label{subsec:extensional-type-theory}
%%%%%%%%%%%%%%%%%%%%%%%%%%%%%%%%%%%%%%%%%%%%%%%%%%%%%%%%%%%%

Extensional approaches impose stronger principles identifying
propositionally equal terms more directly.

This can simplify mathematical reasoning but may complicate computation,
normalization, or decidable type checking.

Different foundational systems make different tradeoffs.

%%%%%%%%%%%%%%%%%%%%%%%%%%%%%%%%%%%%%%%%%%%%%%%%%%%%%%%%%%%%
\subsection{Function Extensionality}
\label{subsec:function-extensionality}
%%%%%%%%%%%%%%%%%%%%%%%%%%%%%%%%%%%%%%%%%%%%%%%%%%%%%%%%%%%%

Mathematically, one often expects:

if

\[
\forall x:A,
\quad
f(x)=g(x),
\]

then

\[
\boxed{
f=g.
}
\]

In type theory this principle is called function extensionality.

It may be derivable, assumed as an axiom, or obtained through stronger
foundational principles depending on the system.

%%%%%%%%%%%%%%%%%%%%%%%%%%%%%%%%%%%%%%%%%%%%%%%%%%%%%%%%%%%%
\subsection{Proof Relevance}
\label{subsec:proof-relevance}
%%%%%%%%%%%%%%%%%%%%%%%%%%%%%%%%%%%%%%%%%%%%%%%%%%%%%%%%%%%%

Under propositions-as-types, proofs are terms.

Therefore two different proofs of the same proposition may be represented
by different terms.

This is called proof relevance.

Some systems or mathematical interpretations intentionally identify more
proofs than others.

%%%%%%%%%%%%%%%%%%%%%%%%%%%%%%%%%%%%%%%%%%%%%%%%%%%%%%%%%%%%
\subsection{Proof Irrelevance}
\label{subsec:proof-irrelevance}
%%%%%%%%%%%%%%%%%%%%%%%%%%%%%%%%%%%%%%%%%%%%%%%%%%%%%%%%%%%%

Proof irrelevance is the principle that all proofs of a proposition are
treated as indistinguishable for some purposes.

Conceptually,

\[
p:P,
\qquad
q:P
\]

may be regarded as equivalent simply because they prove the same
proposition.

Whether this principle holds depends on the type-theoretic setting.

%%%%%%%%%%%%%%%%%%%%%%%%%%%%%%%%%%%%%%%%%%%%%%%%%%%%%%%%%%%%
\subsection{Polymorphism}
\label{subsec:type-theory-polymorphism}
%%%%%%%%%%%%%%%%%%%%%%%%%%%%%%%%%%%%%%%%%%%%%%%%%%%%%%%%%%%%

A polymorphic function can operate uniformly over many types.

For example, the identity function may be represented schematically as

\[
\boxed{
\operatorname{id}
:
\prod_{A:\mathcal{U}}
A\to A.
}
\]

For each type \(A\),

\[
\operatorname{id}_A
\]

acts as the identity on \(A\).

%%%%%%%%%%%%%%%%%%%%%%%%%%%%%%%%%%%%%%%%%%%%%%%%%%%%%%%%%%%%
\subsection{Parametric Polymorphism}
\label{subsec:parametric-polymorphism}
%%%%%%%%%%%%%%%%%%%%%%%%%%%%%%%%%%%%%%%%%%%%%%%%%%%%%%%%%%%%

Parametric polymorphism requires a term to work uniformly for all types.

A polymorphic identity function cannot inspect arbitrary structure inside
\(A\), because no such structure has been supplied.

This uniformity leads to powerful reasoning principles.

%%%%%%%%%%%%%%%%%%%%%%%%%%%%%%%%%%%%%%%%%%%%%%%%%%%%%%%%%%%%
\subsection{System F Preview}
\label{subsec:system-f-preview}
%%%%%%%%%%%%%%%%%%%%%%%%%%%%%%%%%%%%%%%%%%%%%%%%%%%%%%%%%%%%

System F extends simply typed lambda calculus with universal quantification
over types.

A polymorphic identity term has type

\[
\boxed{
\forall A.\,A\to A.
}
\]

System F is important in:

\[
\boxed{
\text{type theory},
}
\]

\[
\boxed{
\text{programming languages},
}
\]

and

\[
\boxed{
\text{proof theory}.
}
\]

%%%%%%%%%%%%%%%%%%%%%%%%%%%%%%%%%%%%%%%%%%%%%%%%%%%%%%%%%%%%
\subsection{Dependent Types and Specifications}
\label{subsec:dependent-types-specifications}
%%%%%%%%%%%%%%%%%%%%%%%%%%%%%%%%%%%%%%%%%%%%%%%%%%%%%%%%%%%%

Dependent types can encode precise program specifications.

For example, one may describe vectors by both element type and length:

\[
\boxed{
\operatorname{Vec}(A,n).
}
\]

Then concatenation can have type

\[
\boxed{
\operatorname{Vec}(A,m)
\to
\operatorname{Vec}(A,n)
\to
\operatorname{Vec}(A,m+n).
}
\]

The type records the length theorem directly.

%%%%%%%%%%%%%%%%%%%%%%%%%%%%%%%%%%%%%%%%%%%%%%%%%%%%%%%%%%%%
\subsection{Worked Example: Length-Indexed Vectors}
\label{subsec:worked-length-indexed-vectors}
%%%%%%%%%%%%%%%%%%%%%%%%%%%%%%%%%%%%%%%%%%%%%%%%%%%%%%%%%%%%

\begin{workedexamplebox}{Dimension in the Type}

Suppose

\[
u:\operatorname{Vec}(\mathbb{R},3)
\]

and

\[
v:\operatorname{Vec}(\mathbb{R},2).
\]

A correctly typed concatenation operation produces

\[
\operatorname{append}(u,v)
:
\operatorname{Vec}(\mathbb{R},5).
\]

\begin{answerbox}
\[
\boxed{
\operatorname{append}(u,v)
:
\operatorname{Vec}(\mathbb{R},3+2).
}
\]
\end{answerbox}

The dimension guarantee is encoded directly in the type.

\end{workedexamplebox}

%%%%%%%%%%%%%%%%%%%%%%%%%%%%%%%%%%%%%%%%%%%%%%%%%%%%%%%%%%%%
\subsection{Refinement Types}
\label{subsec:refinement-types}
%%%%%%%%%%%%%%%%%%%%%%%%%%%%%%%%%%%%%%%%%%%%%%%%%%%%%%%%%%%%

A refinement type describes elements satisfying an additional predicate.

Conceptually,

\[
\boxed{
\{x:A\mid P(x)\}.
}
\]

For example,

\[
\boxed{
\{n:\mathbb{Z}\mid n>0\}
}
\]

represents positive integers.

Different type systems implement refinement types with different logical
strength.

%%%%%%%%%%%%%%%%%%%%%%%%%%%%%%%%%%%%%%%%%%%%%%%%%%%%%%%%%%%%
\subsection{Subtypes}
\label{subsec:type-theory-subtypes}
%%%%%%%%%%%%%%%%%%%%%%%%%%%%%%%%%%%%%%%%%%%%%%%%%%%%%%%%%%%%

A subtype restricts an existing type to terms satisfying a property.

Conceptually,

\[
B
\subseteq
A
\]

can be represented by pairing an element

\[
a:A
\]

with evidence that \(a\) satisfies the defining property of \(B\).

Dependent pair types provide one way to encode this construction.

%%%%%%%%%%%%%%%%%%%%%%%%%%%%%%%%%%%%%%%%%%%%%%%%%%%%%%%%%%%%
\subsection{Quotients}
\label{subsec:type-theory-quotients}
%%%%%%%%%%%%%%%%%%%%%%%%%%%%%%%%%%%%%%%%%%%%%%%%%%%%%%%%%%%%

Mathematics frequently forms quotients by equivalence relations.

Examples include:

\[
\boxed{
\mathbb{Z}/n\mathbb{Z},
}
\]

and quotient spaces in topology.

Representing quotient constructions computationally can be subtle because
one must identify different representatives without losing useful
computational structure.

Some type theories include quotient types explicitly.

%%%%%%%%%%%%%%%%%%%%%%%%%%%%%%%%%%%%%%%%%%%%%%%%%%%%%%%%%%%%
\subsection{Type Checking}
\label{subsec:type-checking}
%%%%%%%%%%%%%%%%%%%%%%%%%%%%%%%%%%%%%%%%%%%%%%%%%%%%%%%%%%%%

Type checking asks:

\[
\boxed{
\text{Given }\Gamma,t,A,
\text{ is }
\Gamma\vdash t:A?
}
\]

In many practical dependent type theories, type checking is designed to be
decidable.

The checker verifies that each term obeys the formal typing rules.

%%%%%%%%%%%%%%%%%%%%%%%%%%%%%%%%%%%%%%%%%%%%%%%%%%%%%%%%%%%%
\subsection{Type Inference}
\label{subsec:type-inference}
%%%%%%%%%%%%%%%%%%%%%%%%%%%%%%%%%%%%%%%%%%%%%%%%%%%%%%%%%%%%

Type inference asks the system to determine a type from an expression.

For example, from

\[
\lambda x.\,x,
\]

one may infer a sufficiently general identity type.

Inference can be easy in some type systems and undecidable or
computationally difficult in more expressive ones.

%%%%%%%%%%%%%%%%%%%%%%%%%%%%%%%%%%%%%%%%%%%%%%%%%%%%%%%%%%%%
\subsection{Type Checking Versus Proof Search}
\label{subsec:type-checking-proof-search}
%%%%%%%%%%%%%%%%%%%%%%%%%%%%%%%%%%%%%%%%%%%%%%%%%%%%%%%%%%%%

Under Curry--Howard:

\[
\boxed{
\text{type checking}
}
\]

corresponds roughly to

\[
\boxed{
\text{proof verification},
}
\]

while

\[
\boxed{
\text{constructing a term of type }P
}
\]

corresponds to

\[
\boxed{
\text{finding a proof of }P.
}
\]

Checking is often substantially easier than search.

%%%%%%%%%%%%%%%%%%%%%%%%%%%%%%%%%%%%%%%%%%%%%%%%%%%%%%%%%%%%
\subsection{Proof Assistants}
\label{subsec:type-theory-proof-assistants}
%%%%%%%%%%%%%%%%%%%%%%%%%%%%%%%%%%%%%%%%%%%%%%%%%%%%%%%%%%%%

Many proof assistants use type theory as their foundational language.

A theorem is represented as a type.

A formal proof is represented as a term inhabiting that type.

The trusted kernel checks a judgment of the form

\[
\boxed{
p:P.
}
\]

Thus machine verification reduces to type checking.

%%%%%%%%%%%%%%%%%%%%%%%%%%%%%%%%%%%%%%%%%%%%%%%%%%%%%%%%%%%%
\subsection{Trusted Kernels}
\label{subsec:type-theory-trusted-kernels}
%%%%%%%%%%%%%%%%%%%%%%%%%%%%%%%%%%%%%%%%%%%%%%%%%%%%%%%%%%%%

A proof assistant may have sophisticated automation for searching for proof
terms.

However, the final proof object is checked by a small kernel implementing
the core type rules.

This creates a separation between:

\[
\boxed{
\text{untrusted proof search}
}
\]

and

\[
\boxed{
\text{trusted proof checking}.
}
\]

%%%%%%%%%%%%%%%%%%%%%%%%%%%%%%%%%%%%%%%%%%%%%%%%%%%%%%%%%%%%
\subsection{Executable Proof Content}
\label{subsec:executable-proof-content}
%%%%%%%%%%%%%%%%%%%%%%%%%%%%%%%%%%%%%%%%%%%%%%%%%%%%%%%%%%%%

Constructive proofs may contain computational data.

For example, a proof of

\[
\boxed{
\forall x:A
\exists y:B
\,
R(x,y)
}
\]

corresponds to a dependent function that takes \(x\) and returns a pair

\[
\boxed{
(y,p)
}
\]

where

\[
y:B
\]

and

\[
p:R(x,y).
\]

Thus the proof itself contains an algorithm for constructing witnesses.

%%%%%%%%%%%%%%%%%%%%%%%%%%%%%%%%%%%%%%%%%%%%%%%%%%%%%%%%%%%%
\subsection{Program Extraction}
\label{subsec:program-extraction}
%%%%%%%%%%%%%%%%%%%%%%%%%%%%%%%%%%%%%%%%%%%%%%%%%%%%%%%%%%%%

In suitable constructive systems, computational content can be extracted
from proofs.

A theorem proving existence of a function may contain enough information to
produce executable code.

This creates a direct bridge between:

\[
\boxed{
\text{mathematical proof}
}
\]

and

\[
\boxed{
\text{verified program}.
}
\]

%%%%%%%%%%%%%%%%%%%%%%%%%%%%%%%%%%%%%%%%%%%%%%%%%%%%%%%%%%%%
\subsection{Type Theory as a Foundation}
\label{subsec:type-theory-foundation}
%%%%%%%%%%%%%%%%%%%%%%%%%%%%%%%%%%%%%%%%%%%%%%%%%%%%%%%%%%%%

Dependent type theory can serve as an alternative foundation for
mathematics.

Instead of saying

\[
x\in A,
\]

one often begins with

\[
\boxed{
x:A.
}
\]

Mathematical structures are represented using types, functions, dependent
pairs, inductive definitions, and identity types.

%%%%%%%%%%%%%%%%%%%%%%%%%%%%%%%%%%%%%%%%%%%%%%%%%%%%%%%%%%%%
\subsection{Set Theory Versus Type Theory}
\label{subsec:set-theory-vs-type-theory}
%%%%%%%%%%%%%%%%%%%%%%%%%%%%%%%%%%%%%%%%%%%%%%%%%%%%%%%%%%%%

Set theory organizes mathematics around membership:

\[
\boxed{
x\in A.
}
\]

Type theory organizes mathematics around typing:

\[
\boxed{
x:A.
}
\]

Set-theoretic foundations emphasize one cumulative universe of sets.

Type-theoretic foundations emphasize structured classes of terms and types.

Both can formalize large portions of modern mathematics.

%%%%%%%%%%%%%%%%%%%%%%%%%%%%%%%%%%%%%%%%%%%%%%%%%%%%%%%%%%%%
\subsection{Advantages of Type-Theoretic Foundations}
\label{subsec:type-theory-foundational-advantages}
%%%%%%%%%%%%%%%%%%%%%%%%%%%%%%%%%%%%%%%%%%%%%%%%%%%%%%%%%%%%

Type theory naturally combines:

\[
\boxed{
\text{logic},
}
\]

\[
\boxed{
\text{computation},
}
\]

and

\[
\boxed{
\text{formal verification}.
}
\]

Well-formedness is enforced syntactically by types.

Constructive proofs can carry executable information.

These features make type theory particularly suitable for mechanized
mathematics.

%%%%%%%%%%%%%%%%%%%%%%%%%%%%%%%%%%%%%%%%%%%%%%%%%%%%%%%%%%%%
\subsection{Homotopy Type Theory Preview}
\label{subsec:homotopy-type-theory-preview}
%%%%%%%%%%%%%%%%%%%%%%%%%%%%%%%%%%%%%%%%%%%%%%%%%%%%%%%%%%%%

Homotopy Type Theory, abbreviated HoTT, interprets types geometrically.

Conceptually:

\[
\boxed{
\text{types}
\leftrightarrow
\text{spaces},
}
\]

\[
\boxed{
\text{terms}
\leftrightarrow
\text{points},
}
\]

and

\[
\boxed{
\text{identity proofs}
\leftrightarrow
\text{paths}.
}
\]

This gives equality higher-dimensional structure.

%%%%%%%%%%%%%%%%%%%%%%%%%%%%%%%%%%%%%%%%%%%%%%%%%%%%%%%%%%%%
\subsection{Higher Identity Structure}
\label{subsec:higher-identity-structure}
%%%%%%%%%%%%%%%%%%%%%%%%%%%%%%%%%%%%%%%%%%%%%%%%%%%%%%%%%%%%

If

\[
p,q:a=_A b,
\]

one may ask whether

\[
p=q
\]

inside another identity type.

Then one may compare proofs of that equality, and continue.

Thus identity types naturally form a hierarchy of higher identifications.

This is the origin of the homotopical interpretation.

%%%%%%%%%%%%%%%%%%%%%%%%%%%%%%%%%%%%%%%%%%%%%%%%%%%%%%%%%%%%
\subsection{Univalence Preview}
\label{subsec:univalence-preview}
%%%%%%%%%%%%%%%%%%%%%%%%%%%%%%%%%%%%%%%%%%%%%%%%%%%%%%%%%%%%

The univalence principle connects equality of types with equivalence of
types.

Informally,

\[
\boxed{
A=B
}
\]

is closely related to

\[
\boxed{
A\simeq B,
}
\]

where

\[
A\simeq B
\]

means that \(A\) and \(B\) are equivalent.

Univalence makes the principle

\[
\boxed{
\text{equivalent structures may be identified}
}
\]

part of the foundational theory itself.

%%%%%%%%%%%%%%%%%%%%%%%%%%%%%%%%%%%%%%%%%%%%%%%%%%%%%%%%%%%%
\subsection{Equivalence of Types}
\label{subsec:type-equivalence}
%%%%%%%%%%%%%%%%%%%%%%%%%%%%%%%%%%%%%%%%%%%%%%%%%%%%%%%%%%%%

A function

\[
f:A\to B
\]

is an equivalence when it has an appropriate inverse up to the equality
notion of the theory.

Thus type equivalence generalizes ordinary bijection or isomorphism.

In homotopy type theory, equivalence is more fundamental than literal
syntactic equality.

%%%%%%%%%%%%%%%%%%%%%%%%%%%%%%%%%%%%%%%%%%%%%%%%%%%%%%%%%%%%
\subsection{Propositions as Special Types}
\label{subsec:propositions-special-types}
%%%%%%%%%%%%%%%%%%%%%%%%%%%%%%%%%%%%%%%%%%%%%%%%%%%%%%%%%%%%

In homotopy-inspired type theory, propositions can be characterized as
types whose inhabitants carry no essential higher distinction.

Informally, a proposition is a type in which any two inhabitants are equal.

Thus propositions form a special class of types within a broader universe
of mathematical structures.

%%%%%%%%%%%%%%%%%%%%%%%%%%%%%%%%%%%%%%%%%%%%%%%%%%%%%%%%%%%%
\subsection{Sets as Special Types}
\label{subsec:sets-special-types}
%%%%%%%%%%%%%%%%%%%%%%%%%%%%%%%%%%%%%%%%%%%%%%%%%%%%%%%%%%%%

A set-like type is one whose identity proofs themselves behave
propositionally.

Thus ordinary sets can be recovered as a certain level of the hierarchy of
types.

This viewpoint embeds set-like mathematics inside a richer theory of
higher structures.

%%%%%%%%%%%%%%%%%%%%%%%%%%%%%%%%%%%%%%%%%%%%%%%%%%%%%%%%%%%%
\subsection{Type Theory and Category Theory}
\label{subsec:type-theory-category-theory}
%%%%%%%%%%%%%%%%%%%%%%%%%%%%%%%%%%%%%%%%%%%%%%%%%%%%%%%%%%%%

Type theory has deep connections with category theory.

Function types resemble exponentials.

Product and sum types resemble categorical products and coproducts.

Dependent types correspond to fibrational and indexed categorical
structures.

The relationship is especially important in categorical logic.

%%%%%%%%%%%%%%%%%%%%%%%%%%%%%%%%%%%%%%%%%%%%%%%%%%%%%%%%%%%%
\subsection{Type Theory and Programming Languages}
\label{subsec:type-theory-programming-languages}
%%%%%%%%%%%%%%%%%%%%%%%%%%%%%%%%%%%%%%%%%%%%%%%%%%%%%%%%%%%%

Programming-language type systems borrow heavily from mathematical type
theory.

Types can enforce properties such as:

\[
\boxed{
\text{input/output structure},
}
\]

\[
\boxed{
\text{data invariants},
}
\]

and

\[
\boxed{
\text{legal program composition}.
}
\]

More expressive type systems can encode stronger specifications.

%%%%%%%%%%%%%%%%%%%%%%%%%%%%%%%%%%%%%%%%%%%%%%%%%%%%%%%%%%%%
\subsection{Type Theory and Formal Verification}
\label{subsec:type-theory-formal-verification}
%%%%%%%%%%%%%%%%%%%%%%%%%%%%%%%%%%%%%%%%%%%%%%%%%%%%%%%%%%%%

A verified program may be represented by a term whose type expresses its
specification.

For example, a sorting function may carry a type expressing that its output
is:

\[
\boxed{
\text{sorted},
}
\]

and

\[
\boxed{
\text{a permutation of the input}.
}
\]

Constructing the program then simultaneously proves correctness.

%%%%%%%%%%%%%%%%%%%%%%%%%%%%%%%%%%%%%%%%%%%%%%%%%%%%%%%%%%%%
\subsection{Type Theory and Constructive Mathematics}
\label{subsec:type-theory-constructive-mathematics}
%%%%%%%%%%%%%%%%%%%%%%%%%%%%%%%%%%%%%%%%%%%%%%%%%%%%%%%%%%%%

Type theory naturally supports constructive mathematics because proofs have
explicit computational content.

A proof of existence gives a witness.

A proof of disjunction identifies a case.

A proof of implication is a transformation.

These ideas are developed further in the next section.

%%%%%%%%%%%%%%%%%%%%%%%%%%%%%%%%%%%%%%%%%%%%%%%%%%%%%%%%%%%%
\subsection{Type-Theoretic Workflow}
\label{subsec:type-theory-workflow}
%%%%%%%%%%%%%%%%%%%%%%%%%%%%%%%%%%%%%%%%%%%%%%%%%%%%%%%%%%%%

When analyzing a type-theoretic construction:

\begin{enumerate}
    \item specify the available primitive types;
    \item state the typing context;
    \item verify formation rules;
    \item identify constructors for the desired type;
    \item identify elimination principles;
    \item state computation rules;
    \item distinguish terms from types;
    \item distinguish typing judgments from propositions about sets;
    \item apply capture-avoiding substitution;
    \item normalize terms when possible;
    \item track definitional equality separately from propositional
          equality;
    \item interpret logical propositions through Curry--Howard when
          appropriate;
    \item use product types for paired data;
    \item use sum types for alternatives;
    \item use function types for transformations;
    \item use \(\Pi\)-types when the output type depends on the input;
    \item use \(\Sigma\)-types when later data depend on earlier data;
    \item use identity types for explicit equality evidence;
    \item use inductive types to define recursively generated data;
    \item use structural recursion for terminating definitions;
    \item use dependent induction for properties of inductive objects;
    \item check universe levels in polymorphic constructions;
    \item distinguish proof search from type checking;
    \item identify which parts of a proof contain executable data;
    \item verify all final proof terms with the formal type checker.
\end{enumerate}

%%%%%%%%%%%%%%%%%%%%%%%%%%%%%%%%%%%%%%%%%%%%%%%%%%%%%%%%%%%%
\subsection{Common Errors}
\label{subsec:type-theory-common-errors}
%%%%%%%%%%%%%%%%%%%%%%%%%%%%%%%%%%%%%%%%%%%%%%%%%%%%%%%%%%%%

\subsubsection{Confusing Membership and Typing}

The judgment

\[
x:A
\]

is not simply another notation for

\[
x\in A.
\]

Typing is primitive in type theory and obeys its own formal rules.

%%%%%%%%%%%%%%%%%%%%%%%%%%%%%%%%%%%%%%%%%%%%%%%%%%%%%%%%%%%%
\subsubsection{Treating Types as Arbitrary Sets}

A type may have computational or higher equality structure not captured by
a naive set interpretation.

Set semantics can be useful, but it is not always definitionally identical
to the type-theoretic viewpoint.

%%%%%%%%%%%%%%%%%%%%%%%%%%%%%%%%%%%%%%%%%%%%%%%%%%%%%%%%%%%%
\subsubsection{Ignoring the Context}

The judgment

\[
t:A
\]

may only be valid under assumptions.

One should write

\[
\boxed{
\Gamma\vdash t:A
}
\]

when context matters.

%%%%%%%%%%%%%%%%%%%%%%%%%%%%%%%%%%%%%%%%%%%%%%%%%%%%%%%%%%%%
\subsubsection{Ignoring Variable Capture}

Substitution inside lambda terms must be capture avoiding.

Bound variables should be renamed before substitution when necessary.

%%%%%%%%%%%%%%%%%%%%%%%%%%%%%%%%%%%%%%%%%%%%%%%%%%%%%%%%%%%%
\subsubsection{Confusing Beta Reduction With Mathematical Equality in General}

Beta reduction is a specific computation rule:

\[
(\lambda x.\,b)(a)
\longrightarrow
b[a/x].
\]

Not every mathematical equality is a beta reduction.

%%%%%%%%%%%%%%%%%%%%%%%%%%%%%%%%%%%%%%%%%%%%%%%%%%%%%%%%%%%%
\subsubsection{Confusing Definitional and Propositional Equality}

Definitional equality is recognized computationally by the formal system.

Propositional equality requires a term of an identity type.

These need not coincide.

%%%%%%%%%%%%%%%%%%%%%%%%%%%%%%%%%%%%%%%%%%%%%%%%%%%%%%%%%%%%
\subsubsection{Assuming Every Type Theory Uses the Same Equality Rules}

Different theories may adopt:

\[
\boxed{
\text{intensional equality},
}
\]

\[
\boxed{
\text{extensional principles},
}
\]

or

\[
\boxed{
\text{univalence}.
}
\]

The exact system must be specified.

%%%%%%%%%%%%%%%%%%%%%%%%%%%%%%%%%%%%%%%%%%%%%%%%%%%%%%%%%%%%
\subsubsection{Assuming Curry--Howard Means Every Programming Language Is a Logic}

Curry--Howard relates particular logical calculi to corresponding typed
computational systems.

The exact correspondence depends on the systems involved.

%%%%%%%%%%%%%%%%%%%%%%%%%%%%%%%%%%%%%%%%%%%%%%%%%%%%%%%%%%%%
\subsubsection{Confusing Product Types and Sum Types}

A term of

\[
A\times B
\]

contains both an \(A\)-term and a \(B\)-term.

A term of

\[
A+B
\]

contains one side together with a tag indicating which side.

%%%%%%%%%%%%%%%%%%%%%%%%%%%%%%%%%%%%%%%%%%%%%%%%%%%%%%%%%%%%
\subsubsection{Confusing Ordinary and Dependent Functions}

In

\[
A\to B,
\]

the codomain is fixed.

In

\[
\prod_{x:A}B(x),
\]

the output type may depend on the input.

%%%%%%%%%%%%%%%%%%%%%%%%%%%%%%%%%%%%%%%%%%%%%%%%%%%%%%%%%%%%
\subsubsection{Confusing Product and Dependent Pair Types}

In

\[
A\times B,
\]

the second type is fixed.

In

\[
\sum_{x:A}B(x),
\]

the second component has type depending on the first component.

%%%%%%%%%%%%%%%%%%%%%%%%%%%%%%%%%%%%%%%%%%%%%%%%%%%%%%%%%%%%
\subsubsection{Treating Existential Proofs as Mere Truth Values}

Constructively, a proof of

\[
\exists x:A\,P(x)
\]

contains both:

\[
\boxed{
\text{a witness}
}
\]

and

\[
\boxed{
\text{evidence for the property}.
}
\]

%%%%%%%%%%%%%%%%%%%%%%%%%%%%%%%%%%%%%%%%%%%%%%%%%%%%%%%%%%%%
\subsubsection{Assuming Every Recursive Definition Is Accepted}

In strongly normalizing dependent type theories, recursive definitions must
often satisfy a structural or well-founded termination condition.

Arbitrary general recursion may destroy normalization.

%%%%%%%%%%%%%%%%%%%%%%%%%%%%%%%%%%%%%%%%%%%%%%%%%%%%%%%%%%%%
\subsubsection{Ignoring Universe Levels}

Expressions that quantify over all types can create universe-size problems.

A universe hierarchy prevents certain circular definitions.

%%%%%%%%%%%%%%%%%%%%%%%%%%%%%%%%%%%%%%%%%%%%%%%%%%%%%%%%%%%%
\subsubsection{Assuming \(\mathcal{U}:\mathcal{U}\) Is Always Safe}

Unrestricted self-typing of universes can create paradoxes.

Modern systems usually impose stratification or other restrictions.

%%%%%%%%%%%%%%%%%%%%%%%%%%%%%%%%%%%%%%%%%%%%%%%%%%%%%%%%%%%%
\subsubsection{Assuming Type Checking and Type Inference Are the Same}

Type checking verifies a supplied type.

Type inference attempts to discover one.

The latter can be significantly harder.

%%%%%%%%%%%%%%%%%%%%%%%%%%%%%%%%%%%%%%%%%%%%%%%%%%%%%%%%%%%%
\subsubsection{Assuming Type Safety Means All Programs Are Correct}

Type safety prevents specified classes of type errors.

It does not automatically guarantee that a program satisfies every desired
mathematical or application-level property.

%%%%%%%%%%%%%%%%%%%%%%%%%%%%%%%%%%%%%%%%%%%%%%%%%%%%%%%%%%%%
\subsubsection{Assuming Strong Normalization Holds in Every Programming Language}

General-purpose programming languages usually permit nontermination.

Strong normalization is a property of restricted typed calculi, not all
computational systems.

%%%%%%%%%%%%%%%%%%%%%%%%%%%%%%%%%%%%%%%%%%%%%%%%%%%%%%%%%%%%
\subsubsection{Confusing Proof Irrelevance With Uniqueness of Proof Syntax}

Proof irrelevance is a mathematical principle about treating proofs as
indistinguishable in a certain sense.

Different syntactic proof expressions may still exist.

%%%%%%%%%%%%%%%%%%%%%%%%%%%%%%%%%%%%%%%%%%%%%%%%%%%%%%%%%%%%
\subsubsection{Assuming Univalence Means Isomorphic Objects Are Literally the Same String}

Univalence concerns equality inside the type-theoretic universe and its
relationship with equivalence.

It is not syntactic identity of expressions.

%%%%%%%%%%%%%%%%%%%%%%%%%%%%%%%%%%%%%%%%%%%%%%%%%%%%%%%%%%%%
\subsection{Applications}
\label{subsec:type-theory-applications}
%%%%%%%%%%%%%%%%%%%%%%%%%%%%%%%%%%%%%%%%%%%%%%%%%%%%%%%%%%%%

\begin{applicationbox}

\textbf{Foundations of Mathematics.}
Dependent type theory provides an alternative foundational framework to
axiomatic set theory.

\medskip

\textbf{Proof Theory.}
The Curry--Howard correspondence identifies formal proofs with typed terms
and normalization with computation.

\medskip

\textbf{Programming Languages.}
Typed lambda calculi provide mathematical models for type systems,
polymorphism, abstraction, and safe program composition.

\medskip

\textbf{Formal Verification.}
Dependent types can encode precise specifications directly in program and
proof types.

\medskip

\textbf{Proof Assistants.}
Machine-checked mathematics can represent theorems as types and proofs as
terms verified by a small trusted kernel.

\medskip

\textbf{Constructive Mathematics.}
Existence proofs produce witnesses, disjunction proofs produce cases, and
universal proofs produce functions.

\medskip

\textbf{Category Theory.}
Products, sums, exponentials, dependent types, and universes have strong
categorical interpretations.

\medskip

\textbf{Homotopy Type Theory.}
Types can be interpreted as spaces, terms as points, and identity proofs as
paths, leading to higher-dimensional foundations.

\end{applicationbox}

%%%%%%%%%%%%%%%%%%%%%%%%%%%%%%%%%%%%%%%%%%%%%%%%%%%%%%%%%%%%
\subsection{Exercises}
\label{subsec:type-theory-exercises}
%%%%%%%%%%%%%%%%%%%%%%%%%%%%%%%%%%%%%%%%%%%%%%%%%%%%%%%%%%%%

\begin{exercise}[1]
Define a typing judgment.
\end{exercise}

\begin{exercise}[1]
Define a typing context.
\end{exercise}

\begin{exercise}[1]
Define a function type.
\end{exercise}

\begin{exercise}[1]
Define a product type.
\end{exercise}

\begin{exercise}[1]
Define a sum type.
\end{exercise}

\begin{exercise}[1]
Define the unit type conceptually.
\end{exercise}

\begin{exercise}[1]
Define the empty type conceptually.
\end{exercise}

\begin{exercise}[1]
State the Curry--Howard correspondence.
\end{exercise}

\begin{exercise}[1]
Define beta reduction.
\end{exercise}

\begin{exercise}[1]
Define a dependent type.
\end{exercise}

\begin{exercise}[2]
Derive

\[
\lambda x.\,x:A\to A.
\]
\end{exercise}

\begin{exercise}[2]
Type the term

\[
\lambda x.\,\lambda y.\,x.
\]
\end{exercise}

\begin{exercise}[2]
Beta reduce

\[
(\lambda x.\,x)(a).
\]
\end{exercise}

\begin{exercise}[2]
Beta reduce

\[
(\lambda x.\,\lambda y.\,x)(a)(b).
\]
\end{exercise}

\begin{exercise}[3]
Explain alpha conversion.
\end{exercise}

\begin{exercise}[3]
Give an example where naive substitution causes variable capture.
\end{exercise}

\begin{exercise}[3]
Perform capture-avoiding substitution in that example.
\end{exercise}

\begin{exercise}[2]
Construct a term of type

\[
A\times B
\]

from terms

\[
a:A
\]

and

\[
b:B.
\]
\end{exercise}

\begin{exercise}[2]
Use projections to construct a function

\[
A\times B\to A.
\]
\end{exercise}

\begin{exercise}[2]
Construct a term of type

\[
A+B
\]

from a term \(a:A\).
\end{exercise}

\begin{exercise}[3]
Explain elimination for a sum type.
\end{exercise}

\begin{exercise}[3]
Explain why the empty type corresponds to falsity.
\end{exercise}

\begin{exercise}[2]
Translate

\[
P\rightarrow Q
\]

into a type.
\end{exercise}

\begin{exercise}[2]
Translate

\[
P\land Q
\]

into a type.
\end{exercise}

\begin{exercise}[2]
Translate

\[
P\lor Q
\]

into a type.
\end{exercise}

\begin{exercise}[2]
Translate

\[
\neg P
\]

into a type.
\end{exercise}

\begin{exercise}[3]
Construct a proof term for

\[
P\land Q\rightarrow Q.
\]
\end{exercise}

\begin{exercise}[3]
Construct a proof term for

\[
P\rightarrow(Q\rightarrow P).
\]
\end{exercise}

\begin{exercise}[3]
Construct a proof term for

\[
P\rightarrow P\lor Q.
\]
\end{exercise}

\begin{exercise}[3]
Explain normalization.
\end{exercise}

\begin{exercise}[3]
Explain strong normalization.
\end{exercise}

\begin{exercise}[3]
Explain preservation.
\end{exercise}

\begin{exercise}[3]
Explain progress.
\end{exercise}

\begin{exercise}[3]
Explain why preservation and progress support type safety.
\end{exercise}

\begin{exercise}[2]
Define a dependent function type.
\end{exercise}

\begin{exercise}[2]
Define a dependent pair type.
\end{exercise}

\begin{exercise}[3]
Show how

\[
A\to B
\]

is a special case of

\[
\prod_{x:A}B(x).
\]
\end{exercise}

\begin{exercise}[3]
Show how

\[
A\times B
\]

is a special case of

\[
\sum_{x:A}B(x).
\]
\end{exercise}

\begin{exercise}[3]
Translate

\[
\forall x:A\,P(x)
\]

into a dependent type.
\end{exercise}

\begin{exercise}[3]
Translate

\[
\exists x:A\,P(x)
\]

into a dependent type.
\end{exercise}

\begin{exercise}[3]
Explain why existential proofs contain witnesses.
\end{exercise}

\begin{exercise}[2]
Define an identity type.
\end{exercise}

\begin{exercise}[2]
State the reflexivity constructor.
\end{exercise}

\begin{exercise}[3]
Explain transport along equality.
\end{exercise}

\begin{exercise}[3]
Distinguish definitional and propositional equality.
\end{exercise}

\begin{exercise}[3]
Give an example of definitional equality by beta computation.
\end{exercise}

\begin{exercise}[2]
Describe the natural numbers as an inductive type.
\end{exercise}

\begin{exercise}[3]
State a recursion principle for natural numbers.
\end{exercise}

\begin{exercise}[3]
State an induction principle for natural numbers.
\end{exercise}

\begin{exercise}[3]
Distinguish recursion from induction.
\end{exercise}

\begin{exercise}[2]
Describe the constructors of a list type.
\end{exercise}

\begin{exercise}[3]
Define list length by structural recursion.
\end{exercise}

\begin{exercise}[3]
Define list append by structural recursion.
\end{exercise}

\begin{exercise}[3]
Describe a binary tree inductively.
\end{exercise}

\begin{exercise}[3]
Explain why structural recursion supports termination.
\end{exercise}

\begin{exercise}[3]
Define canonicity conceptually.
\end{exercise}

\begin{exercise}[3]
Explain how canonicity for \(\mathbb{N}\) gives computational meaning to a
closed natural-number term.
\end{exercise}

\begin{exercise}[3]
Explain how normalization can support consistency.
\end{exercise}

\begin{exercise}[2]
Define a universe conceptually.
\end{exercise}

\begin{exercise}[3]
Explain why universe hierarchies are used.
\end{exercise}

\begin{exercise}[3]
Explain why unrestricted

\[
\mathcal{U}:\mathcal{U}
\]

can be problematic.
\end{exercise}

\begin{exercise}[3]
Distinguish intensional and extensional type theory conceptually.
\end{exercise}

\begin{exercise}[3]
State function extensionality conceptually.
\end{exercise}

\begin{exercise}[3]
Explain proof relevance.
\end{exercise}

\begin{exercise}[3]
Explain proof irrelevance.
\end{exercise}

\begin{exercise}[2]
Define polymorphism conceptually.
\end{exercise}

\begin{exercise}[3]
Give a polymorphic type for the identity function.
\end{exercise}

\begin{exercise}[3]
Explain parametric polymorphism.
\end{exercise}

\begin{exercise}[3]
Describe System F conceptually.
\end{exercise}

\begin{exercise}[3]
Define a length-indexed vector type conceptually.
\end{exercise}

\begin{exercise}[3]
Give a type for vector concatenation preserving lengths.
\end{exercise}

\begin{exercise}[3]
Explain a refinement type.
\end{exercise}

\begin{exercise}[3]
Explain how a subtype can be represented using dependent pairs.
\end{exercise}

\begin{exercise}[3]
Explain why quotient constructions can be subtle computationally.
\end{exercise}

\begin{exercise}[2]
Define type checking.
\end{exercise}

\begin{exercise}[2]
Define type inference.
\end{exercise}

\begin{exercise}[3]
Explain why proof checking corresponds to type checking under
Curry--Howard.
\end{exercise}

\begin{exercise}[3]
Explain the trusted-kernel model of proof assistants.
\end{exercise}

\begin{exercise}[3]
Explain program extraction.
\end{exercise}

\begin{exercise}[3]
Explain how a constructive proof of

\[
\forall x:A\exists y:B\,R(x,y)
\]

contains an algorithm.
\end{exercise}

\begin{exercise}[3]
Compare set-theoretic membership with type-theoretic typing.
\end{exercise}

\begin{exercise}[3]
Explain why dependent type theory can serve as a foundation of mathematics.
\end{exercise}

\begin{exercise}[3]
Explain the homotopy interpretation of types conceptually.
\end{exercise}

\begin{exercise}[3]
Interpret identity proofs as paths.
\end{exercise}

\begin{exercise}[3]
Explain higher identity types conceptually.
\end{exercise}

\begin{exercise}[3]
Explain univalence conceptually.
\end{exercise}

\begin{exercise}[3]
Distinguish type equivalence from syntactic equality.
\end{exercise}

\begin{exercise}[4]
Prove preservation for the simply typed lambda calculus.
\end{exercise}

\begin{exercise}[4]
Prove progress for the simply typed lambda calculus.
\end{exercise}

\begin{exercise}[4]
Study strong normalization of the simply typed lambda calculus.
\end{exercise}

\begin{exercise}[4]
Study the Church--Rosser theorem.
\end{exercise}

\begin{exercise}[4]
Study normalization by evaluation conceptually.
\end{exercise}

\begin{exercise}[4]
Study the Curry--Howard correspondence formally.
\end{exercise}

\begin{exercise}[4]
Develop typing rules for products and sums.
\end{exercise}

\begin{exercise}[4]
Develop dependent \(\Pi\)- and \(\Sigma\)-types formally.
\end{exercise}

\begin{exercise}[4]
Study identity elimination and the \(J\)-rule.
\end{exercise}

\begin{exercise}[4]
Study inductive families.
\end{exercise}

\begin{exercise}[4]
Study dependent pattern matching.
\end{exercise}

\begin{exercise}[4]
Study universe polymorphism.
\end{exercise}

\begin{exercise}[4]
Study impredicative polymorphism and System F.
\end{exercise}

\begin{exercise}[4]
Study refinement types.
\end{exercise}

\begin{exercise}[4]
Study quotient types.
\end{exercise}

\begin{exercise}[4]
Study proof irrelevance and propositional truncation conceptually.
\end{exercise}

\begin{exercise}[4]
Study categorical semantics of type theory.
\end{exercise}

\begin{exercise}[4]
Study Martin-Löf dependent type theory.
\end{exercise}

\begin{exercise}[4]
Study homotopy type theory.
\end{exercise}

\begin{exercise}[4]
Study the univalence axiom.
\end{exercise}

%%%%%%%%%%%%%%%%%%%%%%%%%%%%%%%%%%%%%%%%%%%%%%%%%%%%%%%%%%%%
\subsection{Conceptual Summary}
\label{subsec:type-theory-summary}
%%%%%%%%%%%%%%%%%%%%%%%%%%%%%%%%%%%%%%%%%%%%%%%%%%%%%%%%%%%%

Type theory organizes mathematics through judgments such as

\[
\boxed{
a:A.
}
\]

Types specify which terms are meaningful and how they may be combined.

Function types,

\[
\boxed{
A\to B,
}
\]

represent transformations.

Product types,

\[
\boxed{
A\times B,
}
\]

represent paired data.

Sum types,

\[
\boxed{
A+B,
}
\]

represent alternatives.

Under Curry--Howard,

\[
\boxed{
\text{propositions}
\leftrightarrow
\text{types},
}
\]

and

\[
\boxed{
\text{proofs}
\leftrightarrow
\text{terms}.
}
\]

Dependent types extend this framework:

\[
\boxed{
\prod_{x:A}B(x)
}
\]

represents dependent functions, while

\[
\boxed{
\sum_{x:A}B(x)
}
\]

represents dependent pairs.

Identity types represent equality evidence.

Inductive types provide structured data together with recursion and
induction principles.

Normalization gives computation meaning to proofs, while canonicity ensures
closed data terms reduce to canonical values in suitable systems.

Universes organize types by size.

Modern dependent type theory therefore combines

\[
\boxed{
\text{logic}
+
\text{mathematics}
+
\text{computation}
+
\text{formal verification}.
}
\]

%%%%%%%%%%%%%%%%%%%%%%%%%%%%%%%%%%%%%%%%%%%%%%%%%%%%%%%%%%%%
\subsection{Section Connections}
\label{subsec:type-theory-connections}
%%%%%%%%%%%%%%%%%%%%%%%%%%%%%%%%%%%%%%%%%%%%%%%%%%%%%%%%%%%%

\chapterconnection{
Type Theory develops the proof-as-computation perspective introduced in
Section~\ref{sec:proof-theory}. Computability Theory and Recursion Theory
formalized effective procedures externally through machines and recursive
functions, while Type Theory internalizes computation into the structure of
terms and proofs. The Curry--Howard correspondence identifies propositions
with types and proofs with programs, while dependent types extend this
connection to quantified statements, specifications, equality, and
inductive mathematics. The next section, Constructive Mathematics, develops
the broader mathematical philosophy and methodology in which proofs are
expected to provide explicit witnesses, algorithms, and computational
content rather than relying unrestrictedly on classical principles.
}

%%%%%%%%%%%%%%%%%%%%%%%%%%%%%%%%%%%%%%%%%%%%%%%%%%%%%%%%%%%%
% SECTION SOURCE TRACKING
%%%%%%%%%%%%%%%%%%%%%%%%%%%%%%%%%%%%%%%%%%%%%%%%%%%%%%%%%%%%

%------------------------------------------------------------
% 12.7 Type Theory
%------------------------------------------------------------
%
% Source basis:
%   - Standard simply typed lambda calculus.
%   - Standard Curry--Howard correspondence.
%   - Standard dependent type theory.
%   - Standard inductive types, universes, identity types,
%     proof-assistant foundations, and introductory homotopy type theory.
%
% Used for:
%   - terms and types;
%   - typing judgments;
%   - contexts;
%   - formation, introduction, elimination, and computation rules;
%   - function types;
%   - lambda abstraction;
%   - application;
%   - beta reduction;
%   - simply typed lambda calculus;
%   - product types;
%   - sum types;
%   - unit and empty types;
%   - Curry--Howard correspondence;
%   - implication as function;
%   - conjunction as product;
%   - disjunction as sum;
%   - falsity and negation;
%   - alpha conversion;
%   - substitution;
%   - eta principles;
%   - normalization;
%   - strong normalization;
%   - type safety;
%   - preservation and progress;
%   - dependent types;
%   - Pi types;
%   - Sigma types;
%   - quantified propositions;
%   - identity types;
%   - reflexivity;
%   - equality elimination;
%   - transport;
%   - definitional and propositional equality;
%   - inductive types;
%   - natural-number recursion and induction;
%   - lists and trees;
%   - structural recursion;
%   - canonicity;
%   - normalization and consistency;
%   - universes;
%   - universe hierarchies;
%   - intensional and extensional type theory;
%   - function extensionality;
%   - proof relevance and irrelevance;
%   - polymorphism;
%   - System F preview;
%   - dependent specifications;
%   - length-indexed vectors;
%   - refinement and subtype concepts;
%   - quotient types;
%   - type checking and inference;
%   - proof assistants;
%   - trusted kernels;
%   - program extraction;
%   - type theory as foundation;
%   - relation to set theory;
%   - homotopy type theory preview;
%   - higher identities;
%   - univalence preview;
%   - applications and exercises.
%
% Notes:
%   - Constructive Mathematics is developed next in Section 12.8.
%   - Proof Theory appears in Section 12.4.
%   - Computability Theory appears in Section 12.5.
%   - Recursion Theory appears in Section 12.6.
%   - Philosophy and Foundations follows in Section 12.9.
%
%%%%%%%%%%%%%%%%%%%%%%%%%%%%%%%%%%%%%%%%%%%%%%%%%%%%%%%%%%%%